\chapter{Modélisation du problème}
\label{ch:modelisation}

Ce chapitre formalise le problème du plus court chemin comme une recherche dans
un espace d'états, puis décrit comment une carte OpenStreetMap réelle est
transformée en le graphe pondéré sur lequel opèrent tous les algorithmes.

\section{Le plus court chemin comme recherche dans un espace d'états}
\label{sec:espace-etats}

Un problème de recherche se définit par un espace d'états
$\mathcal{S}$, un état initial $s_0$, une fonction de transition qui associe à
chaque état ses successeurs, un test de but, et un coût de transition. Le plus
court chemin routier s'y inscrit directement :

\begin{bulletlist}
  \item un \textbf{état} est une intersection du réseau ;
  \item les \textbf{transitions} depuis une intersection sont les tronçons de
        route qui en partent ;
  \item le \textbf{coût} d'une transition est la longueur du tronçon (en mètres) ;
  \item le \textbf{test de but} est l'égalité avec l'intersection de destination.
\end{bulletlist}

\paragraph{Graphe routier pondéré.}
Le réseau est modélisé par un graphe orienté pondéré $G = (V, E, w)$ où $V$ est
l'ensemble des intersections, $E \subseteq V \times V$ l'ensemble des tronçons
orientés, et $w : E \to \mathbb{R}_{>0}$ la fonction de poids qui à chaque arc
associe sa longueur. Chaque nœud $v \in V$ porte de plus des coordonnées
géographiques $(\varphi_v, \lambda_v)$ (latitude, longitude).

\medskip

\paragraph{Problème du plus court chemin.}
Étant donnés une source $s \in V$ et une cible $t \in V$, on cherche un chemin
$P = (s = v_0, v_1, \dots, v_k = t)$ avec $(v_{i}, v_{i+1}) \in E$ qui minimise le
coût total
\[
  w(P) \;=\; \sum_{i=0}^{k-1} w(v_i, v_{i+1}).
\]
On note $d(s, t) = \min_P w(P)$ la distance de plus court chemin.

\paragraph{Pourquoi le graphe est orienté.}
Les rues à sens unique imposent que $G$ soit \textbf{orienté} : l'existence de
l'arc $(u, v)$ n'entraîne pas celle de $(v, u)$, et même lorsque les deux
existent la distance n'est pas symétrique en général,
$d(u, v) \neq d(v, u)$. Ce point, souvent négligé, est essentiel : il obligera
l'heuristique ALT (\S\ref{sec:alt}) à stocker \emph{deux} tables de distances par
landmark, l'une sur $G$ et l'autre sur le graphe inverse $G^{\mathsf{R}}$.

\paragraph{Poids strictement positifs.}
Les longueurs de tronçons sont strictement positives. Cette hypothèse
$w > 0$ est ce qui rend Dijkstra correct (un nœud une fois \emph{réglé}, c'est-à-dire
extrait avec sa distance définitive, ne peut plus être amélioré) et ce qui
justifiera l'exactitude de la file à seaux de Dial (\S\ref{sec:bucket}).

\section{D'OpenStreetMap au graphe}
\label{sec:osm}

Le réseau est téléchargé avec la bibliothèque \texttt{osmnx}
\citep{boeing2017osmnx}, qui interroge OpenStreetMap et renvoie le graphe
routier « circulable » d'un lieu. L'extrait par défaut est
\textbf{Paris, France}. On applique ensuite deux traitements indispensables.

\subsection{Composante fortement connexe}

Un graphe routier brut contient des impasses, des zones mal cartographiées et des
nœuds pendants dont on ne peut pas sortir (ou dans lesquels on ne peut pas
entrer) en respectant les sens uniques. Pour garantir que \emph{toute} paire
$(s, t)$ soit joignable dans les deux sens, on ne conserve que la
\textbf{plus grande composante fortement connexe} (SCC) du graphe orienté. Après
ce filtrage, l'extrait de Paris compte :

\begin{center}
\textbf{9\,235 nœuds} \quad et \quad \textbf{17\,886 arcs orientés.}
\end{center}

\subsection{Représentation compacte pour l'échelle}
\label{sec:csr}

Le choix de structure de données est le principal levier d'ingénierie du projet.
Les nœuds sont réindexés en entiers contigus $0, 1, \dots, |V|-1$, et
l'adjacence est stockée dans des tableaux \texttt{numpy} plats de type
\emph{listes d'adjacence compactes} (proches d'un format CSR, \emph{Compressed
Sparse Row}) : un tableau concatène toutes les têtes d'arcs, un tableau parallèle
tous les poids, et un tableau d'index marque, pour chaque nœud, où commencent ses
arcs sortants. Parcourir les voisins d'un nœud revient alors à lire une tranche
contiguë de mémoire, ce qui est à la fois rapide et économe en cache.

\medskip

On maintient de surcroît l'\textbf{adjacence inverse} (les arcs entrants),
nécessaire à la recherche bidirectionnelle (\S\ref{sec:bidir}) et au pré-calcul
d'ALT, ainsi que le \textbf{nom de rue} de chaque arc, réutilisé pour produire un
itinéraire lisible « tourne à gauche / à droite » (\S\ref{sec:multistop}).

\begin{table}[htbp]
  \centering
  \begin{tabular}{@{}ll@{}}
    \toprule
    Propriété & Valeur \\
    \midrule
    Source des données      & OpenStreetMap (via \texttt{osmnx}) \\
    Lieu                    & Paris, France \\
    Filtrage                & plus grande composante fortement connexe \\
    Nombre de nœuds $|V|$   & 9\,235 \\
    Nombre d'arcs $|E|$     & 17\,886 \\
    Poids des arcs          & longueur du tronçon en mètres, $w > 0$ \\
    Orientation             & orienté (sens uniques) \\
    \bottomrule
  \end{tabular}
  \caption{Instance principale utilisée dans tout le rapport.}
  \label{tab:instance}
\end{table}

\section{Géocodage des adresses}
\label{sec:geocodage}

Pour manipuler des trajets réels plutôt que des identifiants de nœuds, une
adresse en toutes lettres (par exemple \emph{« Place du Maréchal de Lattre de
Tassigny, 75016 Paris »}) est géocodée en coordonnées $(\varphi, \lambda)$, puis
projetée sur le \textbf{nœud le plus proche} du graphe. La recherche du nœud le
plus proche utilise la distance orthodromique (\emph{haversine}, définie en
\S\ref{sec:haversine}) entre l'adresse et chaque nœud. Le trajet fil rouge du
rapport, un parcours à plusieurs étapes à travers Paris, est ainsi entièrement
spécifié par des adresses (\S\ref{sec:multistop}).

\section{Métriques évaluées}
\label{sec:metriques}

Conformément à l'énoncé, deux grandeurs mesurent le \emph{travail} d'un
algorithme, indépendamment de la route (identique) qu'il rend :

\begin{bulletlist}
  \item \textbf{Nœuds explorés} : le nombre de nœuds définitivement
        \emph{réglés} (extraits de la file avec leur distance finale). C'est la
        mesure « machine-indépendante » de l'effort, celle qui reflète la qualité
        de l'heuristique.
  \item \textbf{Temps de réponse} : le temps d'exécution moyen par requête,
        mesuré après une requête d'échauffement. Il intègre les \emph{constantes}
        cachées par le décompte de nœuds (coût par nœud de l'heuristique, gestion
        de la file, dépaquetage des raccourcis).
\end{bulletlist}

\noindent
On y ajoute systématiquement le \textbf{contrôle d'optimalité}
(\S\ref{sec:correctness}) : l'erreur relative maximale de longueur de chemin par rapport
à Dijkstra. La distinction entre « nœuds explorés » et « temps » se révélera
cruciale : un algorithme peut explorer moins de nœuds tout en étant plus lent (si
son coût par nœud est élevé), et inversement (Arc Flags, \S\ref{sec:arcflags}).
